\documentclass[10pt,a4paper]{amsart}
\usepackage[margin=19mm]{geometry}
\usepackage{amsmath,amssymb,mathtools}
\usepackage[scr=rsfs,cal=euler]{mathalfa}
\usepackage{xfrac}
\usepackage[unicode,hidelinks,bookmarks=false]{hyperref}
\newcommand{\ie}{i.e.\ }
\newcommand{\cf}{cf.\ }
\newcommand{\resp}{respectively\ }
\newcommand{\Subsection}{\subsection}
\providecommand{\nobreakdash}{\nobreak\hbox{-}}
\setlength{\emergencystretch}{2em}
\multlinegap0pt
\newcommand*{\cg}[1]{\mathcal{C}(#1)}
\usepackage{amssymb}
\usepackage{dsfont}
\usepackage[shortlabels]{enumitem}
\newtheorem{lemma}{Lemma}
\newcommand{\E}{\mathcal{E}}
\newcommand*{\Y}{\mathfrak{Y}}
\newcommand*{\veps}{\varepsilon}
\title{Exhaustion of $\cg{N}$ via rigid expansions}
\author{Jes\'us Hern\'andez Hern\'andez, Cristhian E. Hidber}
\date{Source: arXiv:2603.14736v1 (CC BY 4.0)}
\begin{document}
\maketitle

\noindent\emph{Source and licence.} Exhaustion of $\cg{N}$ via rigid expansions. Version arXiv:2603.14736v1 (CC BY 4.0), distributed by arXiv under the Creative Commons Attribution 4.0 International licence, http://creativecommons.org/licenses/by/4.0/, as recorded in the arXiv licence metadata for this exact version and shown on the abstract page https://arxiv.org/abs/2603.14736v1. Full licence notice: \texttt{licenses/arxiv-2603.14736-CC-BY-4.0.txt}.\par\smallskip

\noindent\textit{Editorial scope.} Counted unit: the article's lemma lemma:nupmij (source0.tex lines 363--365). The article's complete original proof of it is delivered with it (source0.tex lines 366--396, one \texttt{proof} environment of 31 lines). No mathematics is written, repaired or re-derived: the statement and the proof are the article's own lines, in the article's own notation, and no retained line is substituted or re-flowed.  No further source-local context is needed: every cross-reference of every retained line resolves inside the two delivered spans.  One declared adaptation: exactly 1 ancillary strings inside the retained spans are omitted (image-inclusion commands and, where the source repeats a label, the repeated label definitions) and no other line differs from the source; the figure environments, with their placement, captions and labels, are kept, and the package records the omitted strings in fidelity receipts bound to the affected bodies.  No retained line contains a citation, so the wrapper carries no bibliography and no bibliography entry.  The retained lines contain the article's own diagram code, which the wrapper typesets with the standard tikz packages; no image file is delivered and the package declares no asset dependency.  Editorial formatting only, in addition: an AMS \texttt{amsart} preamble that restates the article's own declarations and macros used by the retained lines, namely the environments \texttt{lemma} on separate counters, together with the standard packages those lines need.  The retained environments print in this document as Lemma 1 in order of appearance, whereas the article prints the same items with the numbering of its own full document.  The complete native source of the article as distributed in the arXiv e-print for this exact version is bundled unchanged as \texttt{sources/196429/source0.tex}.
\begin{lemma}\label{lemma:nupmij}
 For all $1 \leq i,j \leq g$ with $|i-j|\geq 2$ (modulo $g$), we have that $\nu^{+}_{i,j}, \nu^{-}_{i,j} \in V(\Y^{6})$.
\end{lemma}

\begin{proof}
Let $1 \leq i,j \leq g$ with $|i-j|\geq 2$ (modulo $g$). In general, we prove that $\nu^{\pm}_{i,j}$ is uniquely determined by a set $A(i,j) \subset V(\Y^{1})$ and an auxiliary curve $\nu$. The former depends only on $i,j$ and the latter depends also of the sign of $\nu^{\pm}_{i,j}$. As such, we describe the set $A(i,j)$, and then use a table to point out what curve is $\nu$.

We define $A(i,j)$ as follows: $$A(i,j) = (\E \setminus \{\veps_{i-2,i},\veps_{i-1,i+1},\veps_{j-2,j},\veps_{j-1,j+1}\}) \cup \{\veps_{i,j-1},\veps_{i-1,j}\} \cup \{\alpha_{k} : 1 \leq k \leq g\} \subset V(\Y^{1}).$$

Before defining $\nu$, we need the following definition to facilitate the exposition.

Let $\mathfrak{i}: \{1, \ldots, g\} \to \mathbb{S}^{1}$ be defined as $k \mapsto \exp(\frac{2k\pi \sqrt{-1}}{g})$. Let also $1 \leq i,j,x,y \leq g$ with $|i-j|\geq 2$ (modulo $g$). We say that $x$ and $y$ are \textit{in different sides from} $i$ and $j$ if $\mathfrak{i}(x)$ and $\mathfrak{i}(y)$ are in different connected components of $\mathbb{S}^{1} \setminus \{\mathfrak{i}(i), \mathfrak{i}(j)\}$; we denote this as $\{x,y\}$ \emph{idsf} $\{i,j\}$.

We describe the curve $\nu$ in Table \ref{table:nu}, along with the rigid expansion to which $\nu^{\pm}_{i,j}$ belongs:

\begin{table}[ht]
 \centering
 \begin{tabular}{|c|c|c|}
  \hline
  $\nu^{\pm}_{i,j}$ & The curve $\nu$ & $m$ for $\nu^{\pm}_{i,j} \in V(\Y^{m})$\\ \hline
  $\nu^{+}_{i,6}$ & $\nu^{-}_{5,7}$ & $2$\\ \hline
  $\nu^{-}_{i,j}$ with $i \neq 6 \neq j$ & $\nu^{+}_{k,6}$ with $\{k,6\}$ \emph{idsf} $\{i,j\}$ & $3$\\ \hline
  $\nu^{+}_{i,j}$ with $\{i,j\} \neq \{5,7\}$ & $\nu^{-}_{k,l}$ with $\{k,l\}$ \emph{idsf} $\{i,j\}$ and $k \neq 6 \neq l$ & $4$\\ \hline
  $\nu^{-}_{i,6}$ & $\nu^{+}_{k,l}$ with $\{k,l\}$ \emph{idsf} $\{i,j\}$ and $\{k,l\}\neq\{5,7\}$ & $5$\\ \hline
  $\nu^{+}_{5,7}$ & $\nu^{-}_{k,l}$ with $\{k,l\}$ \emph{idsf} $\{5,7\}$ & $6$\\ \hline
 \end{tabular}
 \caption{The curves $\nu$ needed to unique determine $\nu^{\pm}_{i,j}$.}\label{table:nu}
\end{table}

See Figures \ref{fig:udet-nupmij}. Therefore, $\nu^{\pm}_{i,j} \in V(\Y^{6})$.
\begin{figure}[ht]
  \centering
  \caption{The curves of the set $A(i,j)$ in blue and the curve $\nu = \nu^{-}_{2,4}$ in red. These curves uniquely determine the curve $\nu^{+}_{3,6}$ in green.}\label{fig:udet-nupmij}
 \end{figure}
\end{proof}

\end{document}
